\chapter{Supply chain con congestione e sostenibilità}
\label{cap:supplychain}

\noindent\textbf{Classe:} LP / NLP convesso \hfill
\textbf{Script:} \texttt{python/lab05\_supplychain.py}

\section{Motivazione gestionale}

I prodotti attraversano una rete di stabilimenti, hub e mercati. Ogni tratta ha un
costo, una capacità e un'impronta di CO$_2$; il costo effettivo cresce rapidamente
quando una tratta si avvicina alla saturazione. Tre domande manageriali:
come instradare i flussi al minimo costo? cosa cambia se penalizziamo la congestione?
quanto deve valere la CO$_2$ perché le rotte ``pulite'' diventino convenienti?



È il problema quotidiano di chi gestisce distribuzione e logistica: scegliere le
rotte, decidere quanto caricare ciascuna tratta, valutare investimenti in capacità
e, sempre più spesso, rendere conto dell'impronta carbonica delle spedizioni. La
struttura matematica --- il flusso su rete --- è tra le più riutilizzabili di tutta
la ricerca operativa.

\textbf{In questo capitolo impareremo a:} (1) formulare un problema di flusso a
costo minimo con vincoli di conservazione; (2) modellare la congestione con termini
quadratici convessi; (3) tracciare una frontiera costo-emissioni con un prezzo
interno della CO$_2$; (4) capire perché le soluzioni degli LP ``saltano'' tra
vertici.

\section{Costruzione guidata del modello}

\paragraph{Il problema a parole.} \emph{Decidiamo} quante unità far viaggiare su
ciascuna tratta della rete. \emph{L'obiettivo} è minimizzare il costo totale di
trasporto (nelle varianti: più congestione o CO$_2$). \emph{I vincoli} impongono che
in ogni nodo i flussi tornino --- gli stabilimenti non spediscono più di quanto
producono, gli hub non trattengono nulla, i mercati ricevono esattamente la loro
domanda --- e che nessuna tratta superi la propria capacità.

\subsection*{Insiemi e dati (input del modello)}
\begin{center}\small
\begin{tabular}{llp{7.4cm}}
\toprule
Simbolo & Tipo & Significato \\
\midrule
$G=(N,A)$ & grafo diretto & nodi $N$ e archi $A \subseteq N \times N$ \\
$b_i$ & $\in \R$ & saldo del nodo $i \in N$: offerta se $> 0$, domanda se $< 0$,
                   transito se $= 0$ \\
$u_{ij}$ & $\in \R_{> 0}$ & capacità dell'arco $(i,j) \in A$ \\
$c_{ij}$ & $\in \R_{\ge 0}$ & costo unitario di trasporto sull'arco $(i,j)$ \\
$e_{ij}$ & $\in \R_{\ge 0}$ & emissioni (kg CO$_2$) per unità trasportata su $(i,j)$ \\
\bottomrule
\end{tabular}
\end{center}

\subsection*{Variabili decisionali}
\begin{center}\small
\begin{tabular}{llp{7.4cm}}
\toprule
$x_{ij}$ & $\ge 0$, continua & flusso inviato sull'arco $(i,j) \in A$ \\
\bottomrule
\end{tabular}
\end{center}

\begin{modello}[Flusso a costo minimo e sue varianti]
\begin{align}
\min\;& \sum_{(i,j) \in A} c_{ij}\, x_{ij} \notag\\
\text{s.t.}\;\;
& \sum_{j:\,(i,j)\in A} x_{ij} \;-\; \sum_{j:\,(j,i)\in A} x_{ji} = b_i && \forall i \in N
   \label{eq:flusso}\\
& 0 \le x_{ij} \le u_{ij} && \forall (i,j) \in A \notag
\end{align}
Varianti sull'obiettivo (tutte con gli stessi vincoli):
\[
\underbrace{\textstyle\sum_{(i,j) \in A} \bigl(c_{ij} x_{ij} + \alpha\, c_{ij} x_{ij}^2 / u_{ij}\bigr)}_{\text{congestione (convessa)}}
\qquad
\underbrace{\textstyle\sum_{(i,j) \in A} (c_{ij} + \tau e_{ij})\, x_{ij}}_{\text{prezzo interno CO$_2$}}
\qquad
\underbrace{\min z \;:\; x_{ij}/u_{ij} \le z \;\; \forall (i,j) \in A}_{\text{minimax utilizzo}}
\]
\end{modello}

\begin{spiegazione}
\textbf{Funzione obiettivo.} Il costo totale di trasporto $\sum_{(i,j)} c_{ij} x_{ij}$:
ogni unità che attraversa l'arco $(i,j)$ paga il suo costo unitario. Le varianti
aggiungono a questo nucleo la congestione o il costo della CO$_2$.

\textbf{Vincolo \eqref{eq:flusso} (conservazione).} In ogni nodo, ciò che esce meno ciò
che entra è pari al saldo $b_i$: positivo negli stabilimenti (producono), negativo nei
mercati (assorbono domanda), nullo negli hub (transito puro). La somma dei $b_i$ su tutti i nodi deve
essere $\ge 0$, altrimenti il problema è inammissibile.

\textbf{Vincolo di capacità ($0 \le x_{ij} \le u_{ij}$).} Nessuna tratta può
trasportare più della propria capacità; è un semplice bound sulle variabili, e il suo
``prezzo ombra'' dirà quanto vale potenziare la tratta.

\textbf{Congestione.} Il termine $\alpha\, c_{ij} x_{ij}^2/u_{ij}$ è convesso e cresce
con l'utilizzo dell'arco: rappresenta straordinari, code ai terminal, premi assicurativi.
Il suo effetto è \emph{ripartire} i flussi su più rotte invece di saturare la più
economica.

\textbf{Prezzo della CO$_2$.} Non serve un vincolo sulle emissioni: basta aggiungerle
al costo con un prezzo interno $\tau$. Variando $\tau$ si traccia l'intera frontiera
costo-emissioni (approccio ``a peso''; l'alternativa ``a vincolo'' impone
$\sum_{(i,j) \in A} e_{ij}\, x_{ij} \le \bar e$ e legge il duale).
\end{spiegazione}

\section{Esempio svolto a mano}

\begin{esempio}[Due rotte in concorrenza]
Un mercato chiede $100$ unità servibili da due rotte: rotta 1 con $c_1 = 2$ \euro{} e
$u_1 = 80$; rotta 2 con $c_2 = 5$ \euro{} e capacità abbondante.

\textbf{LP puro.} Si riempie la rotta economica e si trabocca sull'altra:
$x_1 = 80$, $x_2 = 20$, costo $2 \cdot 80 + 5 \cdot 20 = 260$ \euro.
Prezzo ombra della domanda: servire la 101-esima unità costerebbe $5$ \euro{}
(la rotta 1 è satura); prezzo ombra della capacità $u_1$: $+1$ unità di capacità fa
risparmiare $5 - 2 = 3$ \euro.

\textbf{Con congestione} ($\alpha = 1$): costo rotta 1 $= 2x_1 + 2x_1^2/80$. Il costo
marginale della rotta 1, $2 + 4x_1/80$, va eguagliato a quello della rotta 2
(altrimenti conviene spostare flusso): $2 + x_1/20 = 5 \Rightarrow x_1 = 60$,
$x_2 = 40$. La rotta economica \emph{non viene più saturata}: il modello lascia
margine di sicurezza, proprio come farebbe un pianificatore prudente.
\end{esempio}

\section{Caso di studio}

Rete con 2 stabilimenti (S1, S2), 2 hub (H1, H2) e 4 mercati (M1--M4); 12 archi con
capacità, costi ed emissioni \emph{non} proporzionali ai costi: le tratte economiche
sono ``su strada'' (inquinanti), quelle costose ``su ferro'' (pulite). Dati in
\texttt{dati/supplychain\_archi.csv}; offerta $(260, 240)$, domanda totale $450$.

\begin{lstlisting}[style=python]
m = gp.Model("supply_chain")
x = m.addVars(A, name="x", ub=U)   # ub direttamente sull'arco
m.addConstrs((x.sum(s, "*") <= offerta[s] for s in offerta))
m.addConstrs((x.sum("*", h) == x.sum(h, "*") for h in hub))
v_dom = m.addConstrs((x.sum("*", k) == domanda[k]
                      for k in domanda), name="domanda")
obj = gp.quicksum((c[a] + tau * e[a]) * x[a] for a in A)
if congestione > 0:   # termine quadratico convesso
    obj += gp.quicksum(congestione * c[a] * x[a] * x[a] / U[a]
                       for a in A)
m.setObjective(obj, GRB.MINIMIZE)
\end{lstlisting}

\section{Risultati}

\begin{lstlisting}[style=output]
LP costo minimo : costo 3.385,00 EUR  emissioni 2.730 kg  utilizzo max 100%
  archi saturi: S1->H1, S2->H2, H2->M3
  prezzi ombra della domanda: M1 8,00  M2 9,50  M3 10,00  M4 10,50 EUR/unita'
Congestione a=1: costo 3.702,81 EUR  emissioni 2.530 kg  utilizzo max  90%
Minimax        : utilizzo massimo minimo possibile 58,4%
\end{lstlisting}

\begin{center}
\begin{minipage}{\textwidth}
\centering
% Rete della soluzione: LP a costo minimo (generato da lab05_supplychain.py)
\begin{tikzpicture}[>=stealth,
    nodo/.style={circle, draw=none, text=white, font=\bfseries\small,
                 minimum size=8mm, inner sep=0pt}]
  \draw[rossomattone, line width=2.00pt] (0.00,1.15) -- (3.40,0.92)
    node[midway, above, sloped, font=\tiny, text=black!60] {220};
  \draw[teal, line width=0.47pt] (0.00,1.15) -- (3.40,-0.92)
    node[midway, above, sloped, font=\tiny, text=black!60] {10};
  \draw[black!25, densely dotted, line width=0.4pt] (0.00,-1.15) -- (3.40,0.92);
  \draw[rossomattone, line width=2.00pt] (0.00,-1.15) -- (3.40,-0.92)
    node[midway, above, sloped, font=\tiny, text=black!60] {220};
  \draw[teal, line width=1.27pt] (3.40,0.92) -- (6.80,1.72)
    node[midway, above, sloped, font=\tiny, text=black!60] {120};
  \draw[teal, line width=1.05pt] (3.40,0.92) -- (6.80,0.57)
    node[midway, above, sloped, font=\tiny, text=black!60] {90};
  \draw[teal, line width=0.47pt] (3.40,0.92) -- (6.80,-0.57)
    node[midway, above, sloped, font=\tiny, text=black!60] {10};
  \draw[black!25, densely dotted, line width=0.4pt] (3.40,0.92) -- (6.80,-1.72);
  \draw[black!25, densely dotted, line width=0.4pt] (3.40,-0.92) -- (6.80,1.72);
  \draw[black!25, densely dotted, line width=0.4pt] (3.40,-0.92) -- (6.80,0.57);
  \draw[rossomattone, line width=1.35pt] (3.40,-0.92) -- (6.80,-0.57)
    node[midway, above, sloped, font=\tiny, text=black!60] {130};
  \draw[teal, line width=1.13pt] (3.40,-0.92) -- (6.80,-1.72)
    node[midway, above, sloped, font=\tiny, text=black!60] {100};
  \node[nodo, fill=verde] at (0.00,1.15) {S1};
  \node[nodo, fill=verde] at (0.00,-1.15) {S2};
  \node[nodo, fill=arancio] at (3.40,0.92) {H1};
  \node[nodo, fill=arancio] at (3.40,-0.92) {H2};
  \node[nodo, fill=teal] at (6.80,1.72) {M1};
  \node[nodo, fill=teal] at (6.80,0.57) {M2};
  \node[nodo, fill=teal] at (6.80,-0.57) {M3};
  \node[nodo, fill=teal] at (6.80,-1.72) {M4};
  \node[font=\small\bfseries, text=blunotte] at (3.40,2.30) {LP a costo minimo};
\end{tikzpicture}\\[1.2ex]
% Rete della soluzione: Congestione quadratica (generato da lab05_supplychain.py)
\begin{tikzpicture}[>=stealth,
    nodo/.style={circle, draw=none, text=white, font=\bfseries\small,
                 minimum size=8mm, inner sep=0pt}]
  \draw[teal, line width=1.65pt] (0.00,1.15) -- (3.40,0.92)
    node[midway, above, sloped, font=\tiny, text=black!60] {172};
  \draw[teal, line width=0.73pt] (0.00,1.15) -- (3.40,-0.92)
    node[midway, above, sloped, font=\tiny, text=black!60] {45};
  \draw[teal, line width=0.65pt] (0.00,-1.15) -- (3.40,0.92)
    node[midway, above, sloped, font=\tiny, text=black!60] {35};
  \draw[teal, line width=1.83pt] (0.00,-1.15) -- (3.40,-0.92)
    node[midway, above, sloped, font=\tiny, text=black!60] {197};
  \draw[teal, line width=1.16pt] (3.40,0.92) -- (6.80,1.72)
    node[midway, above, sloped, font=\tiny, text=black!60] {105};
  \draw[teal, line width=0.68pt] (3.40,0.92) -- (6.80,0.57)
    node[midway, above, sloped, font=\tiny, text=black!60] {39};
  \draw[teal, line width=0.69pt] (3.40,0.92) -- (6.80,-0.57)
    node[midway, above, sloped, font=\tiny, text=black!60] {40};
  \draw[teal, line width=0.57pt] (3.40,0.92) -- (6.80,-1.72)
    node[midway, above, sloped, font=\tiny, text=black!60] {23};
  \draw[teal, line width=0.51pt] (3.40,-0.92) -- (6.80,1.72)
    node[midway, above, sloped, font=\tiny, text=black!60] {15};
  \draw[teal, line width=0.77pt] (3.40,-0.92) -- (6.80,0.57)
    node[midway, above, sloped, font=\tiny, text=black!60] {51};
  \draw[teal, line width=1.12pt] (3.40,-0.92) -- (6.80,-0.57)
    node[midway, above, sloped, font=\tiny, text=black!60] {100};
  \draw[teal, line width=0.96pt] (3.40,-0.92) -- (6.80,-1.72)
    node[midway, above, sloped, font=\tiny, text=black!60] {77};
  \node[nodo, fill=verde] at (0.00,1.15) {S1};
  \node[nodo, fill=verde] at (0.00,-1.15) {S2};
  \node[nodo, fill=arancio] at (3.40,0.92) {H1};
  \node[nodo, fill=arancio] at (3.40,-0.92) {H2};
  \node[nodo, fill=teal] at (6.80,1.72) {M1};
  \node[nodo, fill=teal] at (6.80,0.57) {M2};
  \node[nodo, fill=teal] at (6.80,-0.57) {M3};
  \node[nodo, fill=teal] at (6.80,-1.72) {M4};
  \node[font=\small\bfseries, text=blunotte] at (3.40,2.30) {Congestione quadratica};
\end{tikzpicture}
\captionof{figure}{I flussi ottimi sulla rete (stabilimenti S verdi, hub H arancioni,
mercati M teal; lo spessore degli archi è proporzionale al flusso, il rosso indica un
arco saturo, il numero è il flusso in unità). Sopra: LP a costo minimo, che satura
tre archi. Sotto: modello con congestione quadratica, che ripartisce i flussi su più
rotte e non satura alcun arco.}
\end{minipage}
\end{center}

\begin{insight}
I prezzi ombra della domanda ($8$--$10{,}50$ \euro/unità) sono i \emph{costi marginali
di servizio} dei quattro mercati: è il numero da confrontare con il margine commerciale
prima di accettare nuovi ordini, e la base corretta per un prezzo di trasferimento
interno. M4 è il mercato più caro da servire: +31\% rispetto a M1.
\end{insight}

\section{Analisi di sensitività: il prezzo della CO$_2$}

\begin{center}
\begin{tikzpicture}
\begin{axis}[lab, title={Frontiera costo-emissioni al crescere di $\tau$},
  xlabel={emissioni totali (kg CO$_2$)}, ylabel={costo di trasporto (\euro)}]
\addplot+ [mark=*, mark size=1.6pt, nodes near coords={},] table
  [col sep=comma, x=emissioni, y=costo] {\figdat{cap05_frontiera_co2}};
\node [font=\scriptsize, anchor=south east] at (axis cs:2730,3385) {$\tau \le 1$};
\node [font=\scriptsize, anchor=west] at (axis cs:1680,4745) {$\tau \ge 4$};
\end{axis}
\end{tikzpicture}
\captionof{figure}{Frontiera costo-emissioni: ogni punto è la soluzione ottima per un
prezzo interno della CO$_2$ ($\tau$); in basso a destra l'assetto ``su strada''
(economico, 2730 kg), in alto a sinistra l'assetto ``su ferro'' (caro, 1645 kg). La
frontiera procede per salti tra vertici dell'LP.}
\end{center}

\begin{lstlisting}[style=output]
tau = 0.0 ... 1.0 : costo 3385  emissioni 2730   (rotte su strada)
tau = 1.24        : costo 4025  emissioni 2173   (vertice intermedio)
tau = 1.26        : costo 4265  emissioni 1981
tau = 1.5         : costo 4705  emissioni 1661   (assetto "ferro")
tau = 4.0 e oltre : costo 4745  emissioni 1645   (stabile)
\end{lstlisting}

\begin{spiegazione}[Perché la frontiera procede a salti]
Le soluzioni di un LP stanno nei \emph{vertici} del poliedro ammissibile: al crescere
di $\tau$ l'ottimo resta fermo su un vertice, poi \emph{salta} al successivo. Una
bisezione sul modello mostra che fino a $\tau = 1$ non cambia nulla, che tra
$\tau = 1{,}0$ e $\tau = 1{,}5$ la rete attraversa vertici intermedi
(a $\tau = 1{,}24$: 2173 kg) e che il salto principale avviene esattamente a
$\tau = 1{,}25$ \euro/kg. Managerialmente: un prezzo interno della CO$_2$ sotto la
soglia non cambia le rotte; appena sopra, ristruttura la rete in blocco. Il modello
con congestione, strettamente convesso, si sposterebbe invece con continuità.
\end{spiegazione}

\begin{attenzione}
Il minimax ``puro'' ($\min z$) è degenere: qualunque instradamento con utilizzo
$\le z^*$ è ottimo, anche costosissimo. In pratica si usa sempre una combinazione
(costo $+ \rho \cdot z$) o si vincola $z$ e si minimizza il costo --- lo stesso
accorgimento riappare per il peak shaving nel capitolo~\ref{cap:ricaricaev}.
\end{attenzione}

\section{Esercizi}

\begin{esercizio}[verifica]
Verificare per perturbazione il prezzo ombra della domanda di M4 (portarla a 101) e
quello della capacità dell'arco H2$\to$M3.
\end{esercizio}

\begin{esercizio}[guasto di un hub]
Ridurre del 50\% le capacità degli archi entranti e uscenti da H2: quanto costa il
guasto? quali mercati soffrono di più? (Confrontare i prezzi ombra prima/dopo.)
\end{esercizio}

\begin{esercizio}[soglia del ferro]
Riprodurre con una bisezione su $\tau$ la soglia $\tau = 1{,}25$ \euro/kg trovata nel
testo, ed elencare quali archi cambiano flusso attraversandola. Confrontare la soglia
con il prezzo ETS corrente della CO$_2$.
\end{esercizio}

\begin{esercizio}[vincolo sulle emissioni]
Riformulare con vincolo $\sum_{(i,j) \in A} e_{ij}\, x_{ij} \le \bar e$ per $\bar e = 2000$ kg e verificare che
il duale del vincolo coincide con il $\tau$ ``a peso'' che produce le stesse emissioni.
\end{esercizio}

\begin{esercizio}[barriera]
Sostituire il termine di congestione con la barriera
$\alpha\, x_{ij}/(u_{ij} - x_{ij})$ e confrontare: quale formulazione tiene i flussi
più lontani dalla saturazione?
\end{esercizio}
